\chapter{Synthetic Control}
\label{cap:synth}
% ── index ──────────────────────────────────────────────────────
\index{synthetic control|textbf}
\index{synth@\texttt{synth()}|textbf}
\index{RMSPE|textbf}
\index{donors, pool of}
\index{placebo, tests}

% ─────────────────────────────────────────────────────────────────────────────
\section{The Single-Unit Problem}

DiD (ch.~\ref{cap:did}) and PSM (ch.~\ref{cap:psm}) depend on multiple treated
units and comparable controls. In many policy applications, however, only one
unit receives treatment — a country that adopts an economic program, a state
that changes its legislation — and the potentially comparable units are few and
heterogeneous.

\textbf{Synthetic Control} (Abadie, Diamond and Hainmueller, ADH 2010)
constructs an artificial control group — the \emph{synthetic control} — as
a convex combination of donor units that replicates the trajectory of the treated
unit \emph{before} treatment. The central idea:

\[
  \hat{Y}_{1t}^{(0)} = \sum_{j=2}^{J+1} w_j^*\, Y_{jt}
  \qquad
  w_j^* \geq 0,\quad \sum_{j=2}^{J+1} w_j^* = 1
\]

The weights $W^* = (w_2^*,\ldots,w_{J+1}^*)$ minimize the root mean squared
prediction error (\emph{RMSPE}) in the pre-treatment period:

\[
  W^* = \arg\min_{W \in \Delta^J}
  \sum_{t=1}^{T_0-1} \left( Y_{1t} - \sum_{j\geq 2} w_j Y_{jt} \right)^2
\]

where $\Delta^J$ is the unit simplex ($w_j \geq 0$, $\sum w_j = 1$). The treatment
effect in period $t \geq T_0$ is:

\[
  \hat\tau_t = Y_{1t} - \hat{Y}_{1t}^{(0)}
\]

The validity of the method rests on the quality of the pre-treatment fit:
if the synthetic control tracks the treated unit in the pre-treatment period,
the hypothesis that it would have continued to do so in the absence of treatment
is plausible — a non-parametric form of the parallel trends hypothesis.

\begin{nota}
  The convexity restriction ($w_j \geq 0$, $\sum w_j = 1$) avoids
  extrapolation — the synthetic control lies within the support of the donors.
  This distinguishes it from DiD with weights, which can assign negative weights
  to achieve balance.
\end{nota}

% ─────────────────────────────────────────────────────────────────────────────
\section{DGP Verification}

The DGP uses a panel of 10 units $\times$ 20 periods. Unit 1 is treated
from year 11 onward; units 2--10 are donors. The true effect is
$\tau = 3{,}0$ (permanent, from $T_0 = 11$).

\[
  y_{it} = \alpha_i + 0{,}3\,t + 3{,}0\cdot\mathbf{1}[i=1, t\geq 11]
           + \varepsilon_{it}
  \qquad
  \alpha_i \in \{2,3,\ldots,10\},\quad \alpha_1 = 5
\]

The treated unit intercept ($\alpha_1 = 5$) coincides with that of unit 5
(donor), ensuring common support. The innovations $\varepsilon_{it} \sim U(0,1)$
introduce noise that prevents the SC from using only unit 5 as a perfect donor.

\begin{lstlisting}[caption={DGP: 10 × 20 panel, treatment on unit 1 from t=11}]
seed(42)
input df
id
1
end
for i in 1..=8 { df = append(df, df) }
df |> filter(_n <= 200)

generate df unit = (_n - 1) % 10 + 1
generate df year = (_n - (_n - 1) % 10 - 1) / 10 + 1

// intercepts by unit: donors receive alpha = unit; treated alpha = 5
generate df alpha = 0.0
let j = 2
while j <= 10 {
    replace df alpha = j * 1.0 if unit == j
    j = j + 1
}
replace df alpha = 5.0 if unit == 1

generate df d = (unit == 1) * (year >= 11)
generate df e = rnormal()
generate df y = alpha + 0.3 * year + 3.0 * d + e

let m = synth("y", 1, 11, df, id="unit", time="year")
display m
\end{lstlisting}

\subsection*{Panel Construction in Hayashi}

The \lstinline{generate} function operates on the \lstinline{_n} vector (current
row, 1-indexed). For an $N \times T$ panel in long format
(stacked by unit), the formulas are:

\begin{center}
\begin{tabular}{ll}
\toprule
\textbf{Variable} & \textbf{Expression} \\
\midrule
Unit ($1,\ldots,N$)    & \lstinline{(_n - 1) % N + 1} \\
Period ($1,\ldots,T$)  & \lstinline{(_n - (_n-1)%N - 1) / N + 1} \\
\bottomrule
\end{tabular}
\end{center}

\subsection*{Results}

\begin{lstlisting}[language={}, basicstyle=\ttfamily\small,
  frame=none, numbers=none, backgroundcolor=\color{codbg},
  xleftmargin=0pt, xrightmargin=0pt, breaklines=false]
════════════════════════════════════════════════════════════════════
 Synthetic Control  —  Abadie, Diamond, Hainmueller (2010)
════════════════════════════════════════════════════════════════════
 Treated unit : 1   Outcome: y
 T₀ (1st post-treat): 11   Donors: 9   T pre: 10   T post: 10
 RMSPE pre  : 0.3347
 RMSPE post : 2.7056   post/pre ratio: 8.083
────────────────────────────────────────────────────────────────────
 Donor weights (w > 0.001):
   6                         0.4420
   3                         0.2153
   8                         0.1747
   2                         0.1406
   7                         0.0274
────────────────────────────────────────────────────────────────────
  Period         Actual    Synthetic       Effect
 ──────────────────────────────────────────────────
         1        5.8266        6.0335
         2        6.5321        6.1112
         3        6.7458        6.5023
         4        6.6240        6.8732
         5        7.4867        7.2037
         6        7.5689        7.4293
         7        8.0670        7.6433
         8        7.8780        8.0722
         9        8.0202        8.3241
        10        8.0878        8.6960
 *      11       11.4437        8.8167        2.6270
 *      12       12.1888        9.4361        2.7527
 *      13       12.0728        9.8984        2.1745
 *      14       12.9748        9.5753        3.3995
 *      15       12.7653       10.2134        2.5519
 *      16       13.2120       10.2424        2.9696
 *      17       13.3167       10.6701        2.6465
 *      18       13.7205       11.0759        2.6446
 *      19       14.3320       11.5282        2.8039
 *      20       14.0020       11.7106        2.2913
════════════════════════════════════════════════════════════════════
 * post-treatment
\end{lstlisting}

\subsection*{Reading the Results}

\paragraph{Pre-treatment fit.}
The pre-treatment RMSPE is $0{,}33$ for a series whose standard deviation is
on the order of 2.5 — a reasonable fit given the DGP noise ($N=200$, only
10 donors). The post/pre ratio = 8.1 signals that the post-treatment divergence
is \emph{eight times} larger than the pre-treatment fit error, rejecting the
null hypothesis of no effect.

\paragraph{Weights.}
The algorithm distributes weights mainly among units 6, 3, 8 and
2, rather than concentrating all weight on unit 5 (which has $\alpha_5 = 5 =
\alpha_1$). This occurs because the sample noise $\varepsilon_{it}$ makes the
realized trajectory of unit 5 an imprecise linear combination of unit 1
in the pre-treatment period. The SC uses multiple donors to cancel the noise.

\paragraph{Estimated effects.}
The average of the post-treatment effects is:

\[
  \bar{\hat\tau} = \frac{1}{10}\sum_{t=11}^{20}\hat\tau_t = 2{,}69
\]

The true value is $\tau = 3{,}0$; the underestimation of ${\sim}10\%$ reflects
the imprecision of the pre-treatment fit with only 10 donors and 10 pre-periods.

% ─────────────────────────────────────────────────────────────────────────────
\section{Syntax}

\begin{lstlisting}
synth("outcome", treated_id, t0, df, id="unit_col", time="time_col")
synth("outcome", treated_id, t0, df, id="col", time="col", covs=["x1","x2"])
\end{lstlisting}

\begin{center}
\begin{tabular}{lll}
\toprule
\textbf{Argument} & \textbf{Type} & \textbf{Description} \\
\midrule
\texttt{outcome}    & string & outcome variable name \\
\texttt{treated\_id}& num/str& treated unit identifier \\
\texttt{t0}         & number & first post-treatment period \\
\texttt{df}         & DataFrame & panel in long format \\
\texttt{id=}        & option & unit identification column \\
\texttt{time=}      & option & time column \\
\texttt{covs=}      & list   & additional covariates for the fit \\
\bottomrule
\end{tabular}
\end{center}

The \texttt{covs} argument allows including pre-treatment covariate means in
the minimization criterion, in addition to the outcome variable values. Useful
when donors have similar trajectories but distinct structural characteristics.

\begin{lstlisting}[caption={Typical workflow — state policy evaluation}]
load "estados.csv" as df   // id, year, gdp, pop, educ

let m = synth("pib", "SP", 2008, df,
              id="id", time="ano",
              covs=["pop", "educ"])
display m
\end{lstlisting}

\begin{nota}
  The SC does not provide standard asymptotic inference (no standard errors, no
  $p$-values). Inference is conducted via \emph{placebo tests}: the SC is applied
  to each donor unit ``as if'' it were treated and the post/pre ratio is compared.
  If the ratio for the truly treated unit is exceptional in the set of placebos,
  the effect is statistically unusual.
  \hayashi{} displays the post/pre ratio as a descriptive heuristic; placebo tests
  are implemented manually via a loop over units.
\end{nota}

\section{Empirical validation}

The \texttt{synth\_smoking} validation case in \texttt{validation/cases/}
estimates a synthetic control on a simulated panel (10 donors, 1 treated unit,
20 periods). The case compares the average post-treatment effect in \hayashi{}
with a Python reference implementation. Run \lstinline{hay validate} to see the
current status.
